\documentclass[10pt,a4paper]{amsart}
\usepackage[margin=19mm]{geometry}
\usepackage{amsmath,amssymb,mathtools}
\usepackage[scr=rsfs,cal=euler]{mathalfa}
\usepackage{xfrac}
\usepackage[unicode,hidelinks,bookmarks=false]{hyperref}
\newcommand{\ie}{i.e.\ }
\newcommand{\cf}{cf.\ }
\newcommand{\resp}{respectively\ }
\newcommand{\Subsection}{\subsection}
\providecommand{\nobreakdash}{\nobreak\hbox{-}}
\setlength{\emergencystretch}{2em}
\multlinegap0pt
\usepackage{amssymb}
\newtheorem{proposition}{Proposition}
\newcommand{\be}{\begin{equation}}
\newcommand{\bean}{\begin{eqnarray*}}
\newcommand{\ee}{\end{equation}}
\newcommand{\eean}{\end{eqnarray*}}          
\title{The Bishop family of holomorphic discs: regularity and higher index}
\author{Brendan Guilfoyle, Wilhelm Klingenberg}
\date{Source: arXiv:2608.21068v1 (CC BY 4.0)}
\begin{document}
\maketitle

\noindent\emph{Source and licence.} The Bishop family of holomorphic discs: regularity and higher index. Version arXiv:2608.21068v1 (CC BY 4.0), distributed by arXiv under the Creative Commons Attribution 4.0 International licence, http://creativecommons.org/licenses/by/4.0/, as recorded in the arXiv licence metadata for this exact version and shown on the abstract page https://arxiv.org/abs/2608.21068v1. Full licence notice: \texttt{licenses/arxiv-2608.21068-CC-BY-4.0.txt}.\par\smallskip

\noindent\textit{Editorial scope.} Counted unit: the article's proposition p:lin (source0.tex lines 372--378). The article's complete original proof of it is delivered with it (source0.tex lines 379--387, one \texttt{proof} environment of 9 lines). No mathematics is written, repaired or re-derived: the statement and the proof are the article's own lines, in the article's own notation, and no retained line is substituted or re-flowed.  Retained uncounted as the source-local context the proof needs: lines 299--349; lines 351--387.  Nothing was omitted from any retained span, so the package carries no omission record.  No retained line contains a citation, so the wrapper carries no bibliography and no bibliography entry.  No retained line contains a figure environment, an image-inclusion command or drawing code, so no image is delivered and the package declares no asset dependency.  Editorial formatting only, in addition: an AMS \texttt{amsart} preamble that restates the article's own declarations and macros used by the retained lines, namely the environments \texttt{proposition} on separate counters, together with the standard packages those lines need.  The retained environments print in this document as Proposition 1 in order of appearance, whereas the article prints the same items with the numbering of its own full document.  The complete native source of the article as distributed in the arXiv e-print for this exact version is bundled unchanged as \texttt{sources/208223/source0.tex}.
 We arrive at a bundle
\be \label{b}
(p^{-1} A , \Lambda_0^\pm) \stackrel{p}\longrightarrow (A, \partial A).
\ee
This is a real plane bundle over $\overline{A}$ with a real line subbundle $\Lambda_0^\pm$ over the boundary $\partial A$.
A hermitian connection $\nabla$ on the total space induces a decomposition 
\be \label{cs}
 T (p^{-1} A) =TA \oplus VA
\ee   
into horizontal and vertical sub-bundles, where the left hand side stands for the tangent bundle of the total space of (\ref{b}) and $VA$ stands for the fibres of (\ref{b}). Along the base of the bundle (\ref{b}), the complex structure $J$ has the form 
\be \label{J0}
 J|_A \times  0=\bigg(\begin{array}{cc} j & 0\\ 0 &
i\end{array}\bigg).
\ee
Here, $j$ denotes the complex structure on $A$, and $i$ that on the fibre of the tautological bundle at the null dection $A \times  0$. The graph of a section $\varphi\in\Gamma(p^{-1}A,\Lambda_0^\pm)$  of the bundle \eqref{b} is a real surface in the total space whose boundary lies in  $\Lambda_0^\pm$:
\be
{\rm graph} \;\varphi=\{(z,\varphi(z));z \in A\} \subset p^{-1}A,\;\;
\partial\;({\rm graph}\;\varphi) \subset \Lambda_0^\pm. 
\ee
This surface is $J$-holomorphic if its tangent bundle
\[
\mbox{graph}\nabla\varphi= \{(v,\nabla\varphi\cdot v); v\in TA\},
\] 
is $J$-invariant:
\[
\mbox{graph}\nabla\varphi=J\;\mbox{graph}\nabla\varphi
\equiv \{ ((J_{11} + J_{12}\nabla \varphi)v, 
(J_{21} + J_{22}\nabla \varphi)v) ; v\in TA \},
\]
where $J_{ij}$ denote the entries of the matrix of $J \in {\rm Hom}_{\mathbb{R}}(Tp^{-1}A,Tp^{-1}A  )$ with respect to the the decomposition (\ref{cs}), i.e. $J_{11}\in {\rm Hom}_{\mathbb{R}}(TA, TA)$, $J_{12}\in {\rm Hom}_{\mathbb{R}}(VA, TA)$ etc. We now define the differential operator
\be \label{cr}
\bar{\partial}:\Gamma(p^{-1}A, \Lambda_0^{\pm} )\to\Omega^{01}(p^{-1}A),
\;\;\;\bar{\partial}\varphi=\nabla\varphi- K\nabla \varphi  ,
\ee
where  $K\nabla \varphi=(J_{21} + J_{22}\nabla \varphi)\circ (J_{11} + J_{12}\nabla \varphi)^{-1}$. The graphs in $(p^{-1}A,J) $ of solutions of $\bar{\partial} \varphi=0$ are $J$-holomorphic curves on which $J$ induces the orientation that is graphically compatible with that of $(A,j)$. We denote by $\tau$ the $J$-holomorphic involution of the relative bundle \eqref{b} as induced in the preceeding chapter. We have
\[ \tau : (p^{-1}A,\Lambda_0^{\pm}, J) \to (p^{-1}A,\Lambda_0^{\pm}, J).
\]
We are interested in $\tau$-invariant $J$-holomorphic curves of annulus-type $[A \times 0]$ that are sections of $(p^{-1}A,J)$ with boundary values in $\Lambda_0^{\pm}$. That is we seek solutions of
\be \label{cr0}
\bar{\partial}\;\varphi=0,  \;\;\;
\varphi \in \Gamma_\tau(p^{-1}A, \Lambda_0^{\pm}),
\ee
where $\Gamma_\tau$ denotes the $\tau$-invariant sections of \eqref{b}. To this end we linearize the problem as follows. Denote by 
\[
R:= (p^{-1}A|_{\partial A})\cap T\Lambda_0^{\pm}\hookrightarrow p^{-1}A|_{\partial A},
\]
the linearization of the boundary condition, which is a real line sub-bundle of $p^{-1}A|_{\partial A}$. Note that the linearization of $\tau$ along the zero section 
$ O_{p^{-1}A} = A \times 0$ acts on the pair 
\be\label{bdry}
d\tau(O_{p^{-1}A}): (p^{-1}A,R) \to (p^{-1}A,R).
\ee 

The linearization of  $\overline \partial$ in equation \eqref{cr} at the zero section $O$ is defined for  $\psi \in W^{1,p}(p^{-1}A,R)$ by
\be \label{L}
L:W^{1,p}(p^{-1}A,R)\to L^p(\Omega^{01}(p^{-1}A)), \;\;
L(\psi)=d/dt|_{t=0} \bar{\partial}(O_{p^{-1}A} + t\psi). 
\ee
Let $z=(z_0, z_1) \in p^{-1}A$ with $p(z) = z_0$ be the components of a point in the bundle $p^{-1}A$. Then the map
$\tau (z_0, z_1)=(t_0, t_1)$ admits an inverse in the first component $t_0^{-1}(z_0, z_1)$ such that  
$z_0(z_0^{-1}, z_1)=z_0$, and we have 
\bean
 \tau {\rm graph}\; \varphi 
&=& \{ (t_0(z_0, \varphi), t_1(z_0, \varphi))  \; ; z_0 \in A \} \\
&=& {\rm graph}\; t_1 (t_0^{-1},\varphi )\\
&=:& {\rm  graph} \; t \;\varphi.
\eean
Here, last equality is the definition of $t$ as a map of sections.
Now we define the linearization $T$  of $t$ by $\nabla t \varphi=T \nabla \varphi$ with
\be \label{T}
(d  \tau)\; {\rm graph}\nabla \varphi={\rm graph}\; T \nabla \varphi.
\ee
Since $\tau$ is holomorphic, we have $J d \tau=d \tau J $  and $K T= T K$ therefore $\bar{\partial} t\varphi=T \bar{\partial}\varphi$.
\\
\begin{proposition} \label{p:lin}
The operators $L, T$ as in \eqref{L}, \eqref{T} have the following properties:
\begin{enumerate}\item[(a)] $L\psi=\nabla\psi+i
\nabla\psi\circ j + dJ_{21}(O_{p^{-1}A})\cdot \psi \circ j$,
\item[(b)] $L$ is Fredholm,
\item[(c)] L is T-equivariant: $ L\; T =dT\; L$. \end{enumerate} 
\end{proposition}
\begin{proof}
    Part (a) follows from (\ref{J0}) and (\ref{cr}). Note that $dJ_{21}(O_{p^{-1}A})$ is ${\mathbb{R}}$-linear rather than ${\mathbb{C}}$-linear. Part (b) holds since $L$ has the same symbol as the Cauchy-Riemann operator. Part (c) is seen as follows. 
\bean 
L S \psi &=&  d/dr|_{r=0} \;\bar{\partial}\; t \;\varphi_r \\
&=&  d/dr|_{r=0} \; (1-K) \nabla t \varphi_r \\
&=&  (1-d K) \; d T  \;\nabla \psi \\
&=&  d T\; L \psi.
\eean
\end{proof}

\begin{proposition} 
The operators $L, T$ as in \eqref{L}, \eqref{T} have the following properties:
\begin{enumerate}\item[(a)] $L\psi=\nabla\psi+i
\nabla\psi\circ j + dJ_{21}(O_{p^{-1}A})\cdot \psi \circ j$,
\item[(b)] $L$ is Fredholm,
\item[(c)] L is T-equivariant: $ L\; T =dT\; L$. \end{enumerate} 
\end{proposition}

\begin{proof}
    Part (a) follows from (\ref{J0}) and (\ref{cr}). Note that $dJ_{21}(O_{p^{-1}A})$ is ${\mathbb{R}}$-linear rather than ${\mathbb{C}}$-linear. Part (b) holds since $L$ has the same symbol as the Cauchy-Riemann operator. Part (c) is seen as follows. 
\bean 
L S \psi &=&  d/dr|_{r=0} \;\bar{\partial}\; t \;\varphi_r \\
&=&  d/dr|_{r=0} \; (1-K) \nabla t \varphi_r \\
&=&  (1-d K) \; d T  \;\nabla \psi \\
&=&  d T\; L \psi.
\eean
\end{proof}

\end{document}
